\documentclass[11pt,a4paper]{article}
\usepackage[T1]{fontenc}
\usepackage[utf8]{inputenc}
\usepackage{lmodern,microtype}
\usepackage[margin=25mm,headheight=14pt]{geometry}
\usepackage{amsmath,amssymb,amsthm,mathtools}
\usepackage{booktabs,longtable,array,enumitem}
\usepackage[dvipsnames]{xcolor}
\usepackage{xurl,hyperref,fancyhdr}
\definecolor{ink}{HTML}{163B50}
\definecolor{accent}{HTML}{087E8B}
\hypersetup{colorlinks=true,linkcolor=ink,urlcolor=accent,citecolor=ink,
 pdftitle={Smaller Diophantine forms of ZFC consistency},
 pdfauthor={OpenAI-assisted research synthesis for ProveIt}}
\pagestyle{fancy}\fancyhf{}
\fancyhead[L]{\small\textsc{Smaller Diophantine forms of ZFC consistency}}
\fancyhead[R]{\small 22 October 2026}
\fancyfoot[C]{\thepage}
\setlength{\parskip}{4pt}
\setlength{\emergencystretch}{3em}
\setlist{itemsep=3pt,topsep=5pt}
\allowdisplaybreaks[3]
\newtheorem{theorem}{Theorem}[section]
\newtheorem{lemma}[theorem]{Lemma}
\newtheorem{proposition}[theorem]{Proposition}
\theoremstyle{remark}\newtheorem{remark}[theorem]{Remark}
\newcommand{\NN}{\mathbb N}
\newcommand{\ZZ}{\mathbb Z}
\newcommand{\PP}{\mathbb Z_{>0}}
\newcommand{\Con}{\operatorname{Con}}
\newcommand{\ZFC}{\mathrm{ZFC}}
\newcommand{\pc}{\operatorname{pc}}
\newcommand{\pin}{\texttt{ae3007ca4467011e0e051c84acb4a1dd13971b0e}}
\newcommand{\repo}[2]{\href{https://github.com/VladimirReshetnikov/ProveIt/blob/ae3007ca4467011e0e051c84acb4a1dd13971b0e/#1}{#2}}
\newcommand{\collection}{SetTheory/Cardinals/docs/reports/hilbert-tenth-problem/con-zfc-diophantine-forms}
\title{\textbf{Smaller Diophantine forms\\of ZFC consistency}\\[7pt]
\large A twelve-variable polynomial of degree 1081 in 116 operations,\\
with a fivefold smaller coefficient-base exponent}
\author{OpenAI-assisted research synthesis prepared for\\Vladimir Reshetnikov's ProveIt repository}
\date{22 October 2026\\[3pt]\small Comparison snapshot: \texttt{ae3007ca4467}}
\begin{document}
\maketitle
\begin{abstract}
We compare the eight current packages in ProveIt's third collection of
Diophantine formulations of $\Con(\ZFC)$ and combine two of their root
elimination methods with a smaller coefficient encoding. The new ordinary
integer polynomial has twelve strictly positive unknowns, exact total degree
$1081$, and a shared arithmetic circuit with $116$ binary operations.
It improves the inspected twelve-variable degree record $1145$ at $123$
operations and matches the $116$-operation construction of degree $1353$.
Its three-equation presentation has maximum degree $1072$ and $111$
operations under the collection's free-comparison convention.

Separately, mixed-radix clause decoding and external zero tests reduce the
943-tile compiler base from $2^{5^{10}}$ to $2^{5^9}$. All four coefficients
are materialized: their combined raw size falls from $5\,312\,269$ to
$977\,220$ bytes. This tuple also instantiates the existing $67$-operation
polynomial and $59$-operation system, included with unchanged arithmetic.

The proof composes explicit new algebraic and number-theoretic arguments
with the identified native compiler and K215 consistency calibration.
It is ordinary mathematics formalizable in ZFC, supported by reproducible
finite audits, and is not a new end-to-end Lean or Rocq certificate.
No consistency proof, global minimum, or improvement on the existing
nine-natural-variable count is asserted.
\end{abstract}
\clearpage
\begingroup\small\tableofcontents\endgroup
\clearpage

\section{The result and the meaning of simplicity}
\label{sec:result}

The central deliverable is a completely specified polynomial
$\mathcal P_{12}\in\ZZ[J,F,\alpha,z,h,i,T,s,w,\eta,\zeta,r]$.
Its four coefficient ports are bound to four particular integers, not
existential parameters. The full definition appears in
Sections~\ref{sec:formula} and~\ref{sec:coefficients}; executable binary
circuits and the actual coefficient bytes accompany this article.

\begin{theorem}[Fixed-instance equivalence]
\label{thm:main}
For the coefficient tuple defined in Section~\ref{sec:coefficients}, ZFC
proves
\[
 \boxed{\quad
 \bigl(\nexists\mathbf x\in\PP^{12}:\mathcal P_{12}(\mathbf x)=0\bigr)
 \ \longleftrightarrow\ \Con(\ZFC).
 \quad}
\]
The three-equation system and sum-of-squares equation in
Section~\ref{sec:formula} have the same existence predicate.
The polynomial has exact degree $1081$ and a $116$-operation circuit;
the system has maximum residual degree $1072$ and a $111$-operation circuit.
\end{theorem}

Here $\Con(\ZFC)$ means absence of a finite syntactic derivation of falsity
in the effective first-order ZFC presentation used by the K215 endpoint.
The theorem proves an equivalence. It assumes neither side of that
equivalence. The twelve-coordinate domain is strictly positive integers.
For nonnegative unknowns one may write
$\mathcal P_{12}(n_1+1,\ldots,n_{12}+1)=0$; its degree remains $1081$.
The displayed equation is not asserted equivalent over unrestricted
$\ZZ^{12}$.

\subsection{Comparison at an immutable snapshot}

The comparison is frozen at \pin. The current collection has packages
16--23; its earlier waves survive in history and in cited source copies
\cite{catalogue}. The main measured tradeoffs are:

\begin{center}\small
\begin{tabular}{@{}lrrrr@{}}
\toprule
Construction & Unknowns & Equations & Degree & Operations\\
\midrule
18: Jones specialization & $9\NN$ & 1 & $161844\,5^{33}+484$ & 20,531\\
20: previous lowest-degree twelve & $12\PP$ & 1 & 1145 & 123\\
21: previous least-operation twelve & $12\PP$ & 1 & 1353 & 116\\
22: monic twelve & $12\PP$ & 1 & 1529 & 117\\
\textbf{New hybrid} & $\mathbf{12\PP}$ & \textbf{1} & \textbf{1081} & \textbf{116}\\
\textbf{New hybrid system} & $\mathbf{12\PP}$ & \textbf{3} & \textbf{1072} & \textbf{111}\\
New hybrid sum of squares & $12\PP$ & 1 & 2144 & 119\\
Inherited least-operation polynomial & $15\PP$ & 1 & 155 & 67\\
22: inherited least-operation system & $16\PP$ & 7 & 24 & 59\\
Inherited quartic lift & $32\PP$ & 1 & 4 & ---\\
\bottomrule
\end{tabular}
\end{center}

The degree and operation improvement is simultaneous within the reviewed
twelve-variable family: $1081<1145<1353$ and $116<123$.
Relative to package 20 the degree falls by $64$, approximately $5.59\%$;
relative to package 21 it falls by $272$, approximately $20.10\%$.
These statements do not compare unlike domains or imply global optimality.

The $111$-operation system requires an accounting qualification. Our
emitter compares the two final norm terms directly instead of subtracting
them and comparing with zero. The same rewrite reduces package 21's
published $112$-operation system to $111$ operations. Thus the new system
has lower degree at a tied normalized cost. The apparent one-operation
gain is not exclusive to this construction.

An operation means one binary addition, subtraction, or multiplication,
including a square and multiplication by a fixed numeral. Register reuse,
supplied fixed coefficients, and equality comparisons are free. These are
costs of the supplied circuits, not lower bounds on all possible circuits.
We report maximum residual degree for systems and total degree for a
single polynomial.

\subsection{Which presentation to use}

For a compact moderate-degree formula with few coordinates, the new
twelve-variable product is the strongest result here. For arithmetic
operation count alone, the inherited $67$-operation product remains better;
for a small system, retain the $59$-operation presentation. Both are
included with the new, smaller coefficient tuple. If coordinate count
alone is decisive, package 18's nine-natural-variable specialization
remains smaller, at much larger degree and circuit size \cite{jones9}.

The collection also contains $26$-integer and $27$-integer constructions,
including a $27$-integer product of degree $172$ \cite{catalogue}.
Those are different domain tradeoffs. None of the new positive-domain
arguments licenses silently interpreting our twelve coordinates as signed.

\section{A complete polynomial definition}
\label{sec:formula}

Supply exactly the twelve positive integers
\[
 J,F,\alpha,z,h,i,T,s,w,\eta,\zeta,r.
\]
Write the four fixed ports as $\beta,K_0,M_C,M_F$.
The last port is the \emph{shifted} mask
$M_F=M_F^{\rm native}+B-1$. All other symbols below are polynomial
abbreviations, not additional unknowns.

\subsection{The common arithmetic interface}

Define
\begin{align*}
 q&=\beta J+1,& C_w&=q-F-\alpha,& X&=wq,&Y&=sq^3,\\
 E&=XY,&k&=\eta+\zeta,&c&=kY+\eta,&a&=E+Y,\\
 H&=4a+3,&\Delta&=a^2+H=(a+2)^2-1,\\
 \ell&=EkY,&A&=1+\ell(\ell+k),&G&=1+\Delta c^2,\qquad b_0=X+ac.
\end{align*}
The packed rank and outer factors are
\begin{align}
 R&=\bigl((q-1)(q-F)+\alpha\bigr)(q^2-1)
                    +(M_C+qM_F)J,\label{eq:R}\\
 N_k&=k-hE-R,\label{eq:Nk}\\
 N_t&=(K_0+w)C_w+(q-F)-z(q-1).\label{eq:Nt}
\end{align}
For the strong and auxiliary Pell constraints, put
\begin{align*}
 I&=ic^2,&D&=1+\Delta I^2,&Q&=(\Delta I)^2,&d&=Q-1,\\
 v&=cT,&b&=c+RD,&P&=QDv^2,&U&=Qb^2,\\
 e&=P+U-1,& B_{\rm aux}&=2Qvb.
\end{align*}
Our new choice of weights and carrier is
\begin{equation}
 \boxed{M=(kI+T)^2,\qquad L=1+HM,\qquad S=b_0+Hr.}
 \label{eq:weights}
\end{equation}

\subsection{A conceptual norm definition}

For a temporary indeterminate $x$, define
\begin{align*}
 C(x)&=d(x^2+M^2A)-e,\\
 \mathcal F(x)&=
  \bigl(C(x)^2+DB_{\rm aux}^2-4d^2x^2M^2A\bigr)^2
            -4DB_{\rm aux}^2C(x)^2.
\end{align*}
Let $\mathcal F_H(t)=H^8\mathcal F(t/H)$.
The denominator clears because $\deg_x\mathcal F=8$ and all its
coefficients are integers. Define the norm
\begin{equation}
 \mathcal N=\mathcal F_H(S-L\sqrt G)\,
             \mathcal F_H(S+L\sqrt G).
 \label{eq:norm}
\end{equation}
The two conjugate factors make this an integer polynomial. No radical,
division, variable exponent, or quantified auxiliary root occurs in the
actual evaluator supplied next.

\subsection{A division-free scalar evaluator}

The following formulas eliminate all formal radicals explicitly:
\begin{align*}
 n&=(HM)^2A,&V_0&=S^2+L^2G,&V_1&=2SL,\\
 C_0&=d(V_0+n)-H^2e,& C_1&=dV_1,\\
 K_f&=4d^2n,&K_{c0}&=K_fV_0,&K_{c1}&=K_fV_1,&Z&=4H^4PU,\\
 C_{s0}&=C_0^2+GC_1^2,&C_{s1}&=2C_0C_1,\\
 W_0&=C_{s0}-K_{c0}-Z,&W_1&=C_{s1}-K_{c1},\\
 N_0&=W_0^2+GW_1^2-4ZK_{c0},&
 N_1&=2W_0W_1-4ZK_{c1},\\
 \mathcal N&=N_0^2-GN_1^2.
\end{align*}
The fixed coefficient $K_0$ and the intermediate $K_{c0}$ are different
symbols. This distinction is also enforced in the JSON circuits.
The identity used to save operations is
\[
 (C^2+K-Z)^2-4C^2K=(C^2-K-Z)^2-4KZ.
\]
It replaces a general multiplication of two quadratic coefficient pairs
by multiplication of a pair by a scalar. The identity is inherited from
package 21 \cite{norm21}; the distinct inner and outer weights are new.

The recommended single equation is
\begin{equation}
 \boxed{\mathcal P_{12}=N_kN_t(1+\mathcal N)-1=0.}
 \label{eq:main}
\end{equation}
The equivalent three-equation presentation is
\begin{equation}
 \boxed{N_k=1,\qquad N_t=1,\qquad N_0^2=GN_1^2.}
 \label{eq:system}
\end{equation}
The conservative sum-of-squares alternative is
\begin{equation}
 (N_k-1)^2+(N_t-1)^2+\mathcal N^2=0.
 \label{eq:sos}
\end{equation}
Its degree is $2144$ and its circuit uses $119$ operations.

\section{The logical interface and proof dependencies}
\label{sec:dependencies}

Let $K215$ denote the literal machine included as
\texttt{coefficients/k215.tm}. Its working states are $0,\ldots,214$;
state $215$ is halt. The initial tape is blank, with the head in state zero.
The table SHA-256 is
\begin{center}\small\ttfamily
fdbcef0825886035c3091c8537a69535bc96f9e265f0ae45863543d40ae70090.
\end{center}
The pinned endpoint states
\[
 \operatorname{BlankHalts}(K215)
 \ \longleftrightarrow\
 \operatorname{BProv}(\ZFC\mathrm{ax}_s,[],\bot)
 \cite{kendpoint}.
\]
We use the ordinary mathematical calibration mirrored by this endpoint,
not a meta-inference from Lean's soundness to a ZFC proof object.
No endpoint build was rerun for this delivery.

The inherited normalized blank compiler has positive coordinates
\[
 J,F,\alpha,z,f,h,i,T,s,w,\tau,\eta,\zeta,y,\gamma
\]
and the relations
\begin{align*}
 \tau^2&=A,&u&=X+ac+H\gamma,&u^2&=G,\\
 f^2&=1+\Delta i^2c^4=D,&
 Q\bigl(c(Tf-1)-Rf^2\bigr)^2-(Q-1)y^2&=1,\\
 N_k&=1,&N_t&=1.
\end{align*}
Its signed-quotient extension permits the simultaneous replacement
$(f,T)\mapsto(-f,-T)$ during reconstruction; the detailed source
arguments are included in \texttt{upstream/}. This extension, the Pell
rank and half-binomial lemmas, the native compiler, and K215 calibration
are inherited dependencies \cite{blank67,root20,root09,helical}.

The new work has two independent parts: prove that the twelve-variable
norm restores exactly this arithmetic certificate, and prove that the
new coefficient tuple meets the same native compiler interface.
Together they establish
\[
 \exists\mathbf x\in\PP^{12}\;\mathcal P_{12}=0
 \ \longleftrightarrow\ \operatorname{BlankHalts}(K215)
 \ \longleftrightarrow\ \neg\Con(\ZFC).
\]
All statements concern finite strings, natural-number arithmetic, algebraic
numbers, or elementary Pell theory, and the cited constructions are
formalizable in ZFC. The present package does not contain a newly checked
end-to-end formalization of this chain.

\section{The new root-recovery proof}
\label{sec:roots}

\subsection{Unit signs and the noncircular bootstrap}

In the scalar evaluator, $V_1$ and $C_1$ are even, $C_{s1}$ is divisible
by $4$, $K_{c1}$ by $8$, $W_1$ by $4$, and $N_1$ by $8$.
Consequently
\[
 \mathcal N\equiv N_0^2\pmod{64}.
\]
At an integer zero of~\eqref{eq:main}, the three factors
$N_k,N_t,1+\mathcal N$ are units. The value $1+\mathcal N=-1$ would
make a square congruent to $-2$ modulo $4$, which is impossible.
Therefore
\begin{equation}
 \mathcal N=0,\qquad N_k=N_t=\varepsilon\in\{-1,1\}.
 \label{eq:units}
\end{equation}

The inherited weak transport estimate applies before any root recovery
\cite{blank67}. Since $K_0+w\ge2$, the transport factor being a unit
implies $C_w\ge0$. In particular $q-F=C_w+\alpha>0$.
For the coherent compiler tuple its mask bounds then give
\begin{align}
 &q\ge16,\quad X\ge q,\quad Y\ge q^3\ge4096,
 \quad 0<R<q^4-2,\notag\\
 &E\ge q^4>R+2,\quad k=R+\varepsilon+hE>E,\quad k>R,
 \quad kY<c<k(Y+1)<2kY.\label{eq:weak}
\end{align}
The exceptional boundary $C_w=0$ is retained at this stage. It is not
discarded by assuming the final index sign prematurely. Set $A_0=a+2$.
Since $E\ge16Y$ and $Y\ge4096$, we also have $A_0<2E$.

\subsection{The inner weight needs no extra doubling}

Put $P_0=Dv^2+b^2$ and $M_0=kI^2+T^2$.
The estimates~\eqref{eq:weak} give
\begin{equation}
 A>\frac{A_0^2c^2}{16},\qquad A>64A_0^4,\qquad
 D<A_0^2I^2,\qquad b<2kA_0^2I^2.
 \label{eq:innerbounds}
\end{equation}
For the first two, $A>\ell^2$, $A_0c<4EkY=4\ell$, and
$\ell>E^2Y$ suffice. The second follows from
$E^4Y^2>64(2E)^4$. The remaining bounds use $I>1$, $c<I^2$, and $R<k$.

Now $M_0^2\ge4kI^2T^2$ and $M_0^2\ge k^2I^4$. Since $k>64$,
\begin{align*}
 AM_0^2&>\frac{k}{4}A_0^2c^2I^2T^2>16Dv^2,\\
 AM_0^2&>64A_0^4k^2I^4>16b^2.
\end{align*}
Adding these strict inequalities yields $AM_0^2>8P_0$.
Because $d=Q-1\ge Q/2$ and $M=(kI+T)^2\ge M_0$,
\begin{equation}
 \boxed{dAM^2>4QP_0.}\label{eq:inner}
\end{equation}
This estimate uses no root equation. It strengthens the bound used by
package 20 and removes its unnecessary factor of two in the inner weight.

\subsection{All complex branches have a uniform bound}

Factoring $\mathcal F$ first with $f^2=D$ and then with $\tau^2=A$
gives exactly the coherent branches
\begin{equation}
 x=M\tau+y,\qquad \tau^2=A,\quad f^2=D,\quad
 dy^2=e-B_{\rm aux}f=Q(vf-b)^2-1.
 \label{eq:branches}
\end{equation}
This factorization is valid over an algebraic closure, including repeated
roots. It does not introduce assumptions on independent square classes.
Every branch, including a possible imaginary $y$, satisfies
\begin{equation}
 |y|^2\le\frac{2QP_0+1}{d}<\frac{4QP_0}{d}<AM^2.
 \label{eq:ybound}
\end{equation}
It also satisfies the useful separate estimate
\[
 |y|<2(v\sqrt D+b)
       <8k^2YIT+10k^3I^2.
\]
Indeed $\sqrt D<A_0I$, $c<2kY$, $A_0<2E<2k$, and
$b<2kY+4k^3I^2<5k^3I^2$ prove the last inequality.

\begin{lemma}[The main radicand has a large square-root gap]
\label{lem:gap}
Under~\eqref{eq:weak}, $\sqrt G-\sqrt A>kY^2$.
\end{lemma}
\begin{proof}
Since $\Delta>(a+1)^2$ and $c>kY$,
\[
 \sqrt G>(a+1)c>(E+Y+1)kY.
\]
On the other hand,
$(\ell+k/2)^2-A=k^2/4-1>0$, hence
$\sqrt A<\ell+k/2=EkY+k/2$. Subtraction gives
$\sqrt G-\sqrt A>k(Y^2+Y-1/2)>kY^2$.
\end{proof}

Because $M=k^2I^2+2kIT+T^2$ and $Y\ge4096$, this gap gives
\begin{align*}
 M(\sqrt G-\sqrt A)
 &>k^3Y^2I^2+2k^2Y^2IT\\
 &>10k^3I^2+8k^2YIT>|y|.
\end{align*}
Thus every root of $\mathcal F$ obeys
\begin{equation}
 \boxed{|x|<M\sqrt G<\frac{L\sqrt G}{H}.}
 \label{eq:outerbound}
\end{equation}
This is the key saving: the inner weight $M$ is separated from the outer
congruence weight $L=1+HM$. One does not need the extra large outer
multiplier used in package 20, or the repeated outer weight used in 21.

\subsection{The outer norm forces the main root to be integral}

Suppose $\mathcal N=0$. If $G$ were nonsquare, expressing either
factor in~\eqref{eq:norm} as $E_0+E_1\sqrt G$, with integers $E_0,E_1$,
would show that both coefficients vanish: otherwise $G$ is a rational
square and therefore an integer square. Both conjugate factors would
then vanish. The two corresponding roots
$(S-L\sqrt G)/H$ and $(S+L\sqrt G)/H$ differ by
$2L\sqrt G/H$, contradicting~\eqref{eq:outerbound}.

Therefore an integer $u$ with $u^2=G$ and a rational
$x=(S-Lu)/H$ satisfy $\mathcal F(x)=0$.
The root $u$ is positive. If it were negative,
$S=Lu+Hx<0$ by~\eqref{eq:outerbound}, whereas $S=b_0+Hr>0$.

\subsection{Rational reconstruction and Pell saturation}

For real $x$, inequality~\eqref{eq:inner} and
$B_{\rm aux}\sqrt D\le QP_0$ give
\[
 C(x)>3QP_0+1>B_{\rm aux}\sqrt D>0.
\]
Moreover $\mathcal F(0)=(C(0)^2-DB_{\rm aux}^2)^2>0$.
Define
\begin{align*}
 V(x)&=C(x)^2+DB_{\rm aux}^2-4d^2x^2M^2A,\\
 f&=-\frac{V(x)}{2C(x)B_{\rm aux}},&
 \tau&=\frac{C(x)+B_{\rm aux}f}{2dxM},&
 y&=x-M\tau.
\end{align*}
All denominators are nonzero. The equation $\mathcal F(x)=0$ gives
$f^2=D$, and the ensuing identity
$(C+B_{\rm aux}f)^2=4d^2x^2M^2A$ gives $\tau^2=A$.
Both roots are rational numbers with integer squares, so are integers.
The remaining rational $y$ satisfies
\begin{equation}
 Q(vf-b)^2-(Q-1)y^2=1.\label{eq:saturationinstance}
\end{equation}

\begin{lemma}[Rational auxiliary Pell saturation]
\label{lem:saturation}
If integers $m\ge2,V$ and rational $y$ satisfy
$m^2V^2-(m^2-1)y^2=1$, then $y$ is an integer.
\end{lemma}
\begin{proof}
Write $y=a/b$ in lowest terms with $b>0$. The equation implies
$b^2\mid m^2-1$. Put $D_0=(m^2-1)/b^2$, a positive nonsquare integer.
There are integral Pell solutions
\[
 m^2-D_0b^2=1,\qquad (m|V|)^2-D_0a^2=1.
\]
Neither second coordinate is zero. Let
$\epsilon=x_1+y_1\sqrt{D_0}>1$ be the fundamental integral Pell unit,
and write $\epsilon^n=\chi_n+\psi_n\sqrt{D_0}$.
For some $j,l\ge1$, $m=\chi_j$, $b=\psi_j$,
$m|V|=\chi_l$, and $|a|=\psi_l$.

The identity $\chi_j\mid\chi_l$ forces $l$ to be an odd multiple of $j$.
To see this, $\epsilon^{2j}\equiv-1$ in the quadratic ring modulo
$\chi_j$ gives $\chi_{n+2j}\equiv-\chi_n$ and
$\chi_{2j-n}\equiv-\chi_n$. Reduce $l$ modulo $2j$ and reflect into
$0\le n\le j$. For $n<j$, $0<\chi_n<\chi_j$, so divisibility is
impossible; hence the reduced index is $j$.
The Pell multiplication formulas now give $\psi_j\mid\psi_l$.
Thus $b\mid a$, and coprimality forces $b=1$.
\end{proof}

Apply the lemma to $m=\Delta I$ and $V=vf-b$ in
\eqref{eq:saturationinstance}. It yields $y\in\ZZ$, and then
$x=M\tau+y\in\ZZ$. The special Pell shape is essential; a general
rational solution of $dy^2=e-Bf$ need not be integral.

\subsection{Recover the omitted congruence and positive coordinate}

Since $L\equiv1\pmod H$ and $S=b_0+Hr=Lu+Hx$,
$u\equiv b_0\pmod H$. Therefore
\[
 \gamma=\frac{u-b_0}{H}\in\ZZ.
\]
It is positive. Write $\psi_p(n)$ for the coefficient of
$\sqrt{p^2-1}$ in $(p+\sqrt{p^2-1})^n$.
The restored first Pell equation implies
$k=2\psi_{2XY^2+1}(n)$ with $n\ge1$. Since $k>E\ge16^4$,
$n\ne1$, so $k\ge4(2XY^2+1)>X$ and $c>k>X$.
Because $u>0$, $u^2=1+\Delta c^2$, and $\Delta>(a+1)^2$,
\[
 \gamma H=u-ac-X>c-X>0.
\]
Taking absolute values of $\tau$ and $y$ is harmless. If $f<0$, change
the signs of $f,T$ together. The inherited signed-quotient normalized
compiler then applies, forces $\varepsilon=1$, and recovers a blank
halting tableau. This use of the signed-quotient result is explicit;
no signed-$\gamma$ theorem is being assumed.

Conversely, take a canonical positive normalized compiler witness, with
$u=b_0+H\gamma$. The choice
\begin{equation}
 r=\gamma+M(u+\tau)+y>0\label{eq:extension}
\end{equation}
gives $S=Lu+H(M\tau+y)$ and makes a factor of $\mathcal N$ vanish.
Together with $N_k=N_t=1$, it proves completeness. This finishes the new
arithmetic part of Theorem~\ref{thm:main}.

\section{Exact degrees and simpler-proof alternatives}
\label{sec:degree}

Give each supplied unknown degree one and each fixed port degree zero.
The leading-degree ledger is
\begin{center}\small
\begin{tabular}{@{}lr@{\qquad}lr@{\qquad}lr@{}}
\toprule
Expression & Degree & Expression & Degree & Expression & Degree\\
\midrule
$c$ & 5 & $a,H$ & 6 & $\Delta$ & 12\\
$\ell,I$ & 11 & $A,G$ & 22 & $D$ & 34\\
$Q,d$ & 46 & $b$ & 38 & $P$ & 92\\
$U,e$ & 122 & $M$ & 24 & $L,HM$ & 30\\
$S$ & 11 & $n,V_0$ & 82 & $V_1$ & 41\\
$C_0$ & 134 & $C_1$ & 87 & $Z$ & 238\\
$K_{c0}$ & 256 & $K_{c1}$ & 215 & $W_0$ & 268\\
$W_1$ & 221 & $N_0$ & 536 & $N_1$ & 489\\
\bottomrule
\end{tabular}
\end{center}
The unique top contribution to $C_0$ is $-H^2e$.
It remains the unique largest contribution through the subsequent
squares. In particular $N_0^2$ has degree $1072$, while
$GN_1^2$ has degree only $1000$.
Writing $q_0=\beta J$ and $k_0=\eta+\zeta$ gives the uniform certificate
\begin{equation}
 \mathcal N_{\rm top}
 =2^{32}i^{48}k_0^{96}w^{80}s^{176}q_0^{656}(q_0-F)^{16}.
 \label{eq:leader}
\end{equation}
This is nonzero for every valid $\beta>0$. The leading forms of $N_k,N_t$
have degrees $7,2$, respectively, so $\deg\mathcal P_{12}=1081$.
More explicitly, the coefficient of
\[
 J^{677}hz\,i^{48}\eta^{96}w^{81}s^{177}
\]
is $2^{32}\beta^{677}>0$. Thus exact degree is established for the actual
coefficient instance, without relying on small diagnostic replacements.

There is a useful proof-length tradeoff. Replace the carrier by
$S=Lb_0+Hr$, leaving $M=(kI+T)^2$ and $L=1+HM$ unchanged.
This costs one extra multiplication and retains degree $1081$.
The inner bound gives $|x|<2M\sqrt A$, while both $\sqrt G>\sqrt A$
and $b_0>\sqrt A$. For the second inequality,
\[
 a^2k^2Y^2-A=k^2\bigl[EY(2Y^2-1)+Y^4\bigr]-1>0,
\]
and $b_0>akY$. Consequently
$(S+L\sqrt G)/H>2M\sqrt A$, excluding that conjugate directly.
This avoids the outer gap estimates. Its product and system cost
$117$ and $112$ operations. A further variant using $M=kI^2+T^2$
costs $118$ and $113$. All three variants are emitted and audited.

\section{A smaller complete coefficient tuple}
\label{sec:coefficients}

The new adapter retains package 19's $943$ Wang tiles, occupied binary
edge labels, native support, and anchor ruler \cite{wang943,compact943}.
It combines three changes: decode the selector count before other fields;
place mask-population padding in guaranteed zero cells; and shorten the
spatial and clause shifts using exact support checks.

\subsection{The fixed tile language}

A tile is $(\mathrm{west},\mathrm{north},\mathrm{east},\mathrm{south})$.
Use vertical labels $0,1$ for headless bits,
$\operatorname{head}(q,b)=2(q+1)+b$ for $0\le q\le215$, and
$L_0=434,Q_0=435,R_0=436,V_0=437,H_0=438$ for initialization and boundaries.
Horizontal labels are $O=0,HB=1,PL=2,PQ=3,PR=4,BAD=5$ and directed
signals. Let $S_R,S_L$ be the sorted sets of states targeted by right-
and left-moving working actions. They have sizes $199$ and $45$.
Set $\operatorname{right}(q)=6+\operatorname{index}_{S_R}(q)$ and
$\operatorname{left}(q)=205+\operatorname{index}_{S_L}(q)$.

The finite language is specified by the following table. Each range is
finite, and the builder reconstructs it without running the machine.
\begin{center}\small
\begin{tabular}{@{}>{\raggedright\arraybackslash}p{.79\linewidth}r@{}}
\toprule
Tile family & Count\\
\midrule
For each working action $(q,b)\mapsto(v,R,t)$ use
$(O,\operatorname{head}(q,b),\operatorname{right}(t),v)$;
for a left action use $(\operatorname{left}(t),\operatorname{head}(q,b),O,v)$ & 430\\
Halt loops $(O,\operatorname{head}(215,b),O,\operatorname{head}(215,b))$ & 2\\
Headless stays $(O,b,O,b)$ & 2\\
Receive-right $(\operatorname{right}(q),b,O,\operatorname{head}(q,b))$;
receive-left $(O,b,\operatorname{left}(q),\operatorname{head}(q,b))$ & 488\\
Finish $(O,c,O,H_0)$, $c\in\{0,1,\operatorname{head}(215,0),\operatorname{head}(215,1)\}$ & 4\\
Initialization $(PL,L_0,PL,0)$, $(PL,Q_0,PQ,1)$,
$(PQ,R_0,PR,\operatorname{head}(93,0))$, $(PR,R_0,PR,0)$ & 4\\
Reset $(HB,H_0,HB,p)$, $p\in\{L_0,Q_0,R_0\}$ & 3\\
Seam $(u,V_0,v,V_0)$, $u\in\{O,HB,PR\}$, $v\in\{O,HB,PL\}$ & 9\\
Forbidden $(BAD,0,O,0)$ & 1\\
\midrule
Total & 943\\
\bottomrule
\end{tabular}
\end{center}
The start tile $(PL,Q_0,PQ,1)$ is selector zero and the forbidden tile
is selector one; all remaining tiles follow lexicographically.
The actual first machine action is $(0,0)\mapsto(1,R,93)$.
No east color is $BAD$, so the forbidden tile cannot occur in a global
matching tiling.

The inherited language proof shows that a marked periodic tiling exists
exactly when blank K215 halts. Initialization phases force finite vertical
boundaries around the marked head, directed signals implement each exact
machine step, and a later reset row requires a halted head. Conversely a
finite halting run admits independently padded widths and heights and
matching helical seams. The adapter changes no tile or machine action.

\subsection{Literal coefficient definitions}

Put
\[
 \rho=2^{125},\qquad L_{\rm cell}=15625,\qquad
 d_{\rm base}=125\cdot15625=5^9,
 \qquad B=\rho^{L_{\rm cell}}.
\]
For $0\le j<10$, let
\[
 e_0(j)=1023+j,\quad e_1(j)=1013+j,\quad
 e_2(j)=1003+j,\quad e_3(j)=983+j.
\]
The native support is
\[
 \mathcal E=\{0,\ldots,942\}\cup
 \{e_t(j):0\le t<4,\ 0\le j<10\}
 \cup\{1063,1064,1095,1096,2195,3296\}.
\]
It has $989$ elements. Let $w_j=2^{10+2j}$ for $0\le j<44$.
If tile $i$ has labels $x_{it}$, encode the occupied label $x_{it}+512$
in ten bits, and define
\begin{align*}
 c_i&=1+\sum_{t=0}^3\sum_{j=0}^9
          \operatorname{bit}_j(x_{it}+512)w_{10t+j}
             +\sum_{j=40}^{43}w_j,\quad 0\le i<943,\\
 c_{e_t(j)}&=w_{10t+j},\\
 (c_{1095},c_{1096},c_{2195},c_{3296})&=(w_{40},w_{41},w_{42},w_{43}),\\
 c_{1063}&=c_{1064}=0,\qquad
 \mu=1022+\sum_{j=0}^{43}w_j.
\end{align*}
Bit extraction here is only part of the bounded definition of fixed
coefficients from the finite tile list. It is not an operation in the
Diophantine polynomial.

The external padding mask is
\[
 Z_{\rm pad}=(2^{123}-1)\sum_{e=9000}^{9005}\rho^e
                  +(2^{122}-1)\rho^{9006}.
\]
Define
\begin{align}
 D_C={}&\rho^{30}+\rho^{62}+\rho^{1100}+\rho^{2200}
       +\sum_{e\in\mathcal E}c_e
          \bigl(\rho^{4393-e}+\rho^{5495-e}\bigr),\notag\\
 D_R={}&\rho^{42},\notag\\
 M_C={}&B-1-\sum_{\substack{e\in\mathcal E\\e\ne1}}\rho^e
                       -2\rho^{1063},\notag\\
 M_F^{\rm native}={}&4+\mu(\rho^{4393}+\rho^{5495})
  +\sum_{j=0}^9(\rho^{1013+j}+\rho^{1065+j})
  +\rho^{2195}+\rho^{3296}+Z_{\rm pad}.\label{eq:constants}
\end{align}
The complete tuple used throughout this package is
\begin{equation}
 \boxed{\beta=B-1,\qquad K_0=D_C+BD_R,\qquad
          M_C,\qquad M_F=M_F^{\rm native}+B-1.}
 \label{eq:ports}
\end{equation}
Every exponent and summation range is fixed. There is no unbounded
coefficient search and no dependence on whether a refutation is found.

\section{Why the smaller encoding preserves the compiler}
\label{sec:adapterproof}

\subsection{Mixed-radix decoding must be done in order}

At a clean clause target, let $n$ be the number of selected tiles.
Native selector digits are Boolean, so $0\le n\le943<1024$.
All other contributions to the clause sum are divisible by $1024$;
therefore its low ten bits are exactly $n$, even if higher fields have
not yet been decoded. Disjointness from the low mask $1022$ forces
$n\in\{0,1\}$.

Only now consider the upper fields. Each is the sum of at most one
selected tile bit and one copy bit, or of $n$ and one anchor bit.
Its value lies in $\{0,1,2\}$, so radix-four fields are carry-free.
The mask tests their low bits, allowing precisely $0$ and $2$.
Thus an empty cell has no copies or anchors, and an occupied cell has
exactly the selected tile's copy bits and all four anchors.

Conversely, every valid tile encoding and the empty encoding passes.
The old uniform ten-bit spacing is unnecessary after the selector count
has been established. The new clause mask has $9+44=53$ set bits.
Assuming the upper fields carry-free before decoding $n$ would be an
error; the ordered argument above is the substantive mixed-radix change.

\subsection{Exact geometry replaces excessive spacing}

Use clause targets $T_1=4393,T_2=5495$, spatial shift $H_c=42$,
and bare shifts $30,62,1100,2200$.
The required target identities, for $0\le j<10$, are:
\begin{center}\small
\begin{tabular}{@{}lccc@{}}
\toprule
Target & $D_C$ source & $D_R$ source & Temporal source\\
\midrule
$1013+j$ & $983+j$ & none & $1013+j$\\
$1065+j$ & $1003+j$ & $1023+j$ & none\\
$2195$ & $1095$ & none & $2195$\\
$3296$ & $1096$ & none & $3296$\\
\bottomrule
\end{tabular}
\end{center}
For every $e\in\mathcal E$, these are exact native-basis coefficient
identities, including both reflected clause bands. At each clause target
$T$, $[\rho^{T-e}]D_C=c_e$ and $e+42\ne T$.
The checker exhausts all $69\,230$ such identities; linearity extends
them to every native assignment, including the allowed dummy digits.

The two finite support intersections are
\begin{align*}
 (T_1-\mathcal E)\cap(T_2-\mathcal E)&=\varnothing
                         \pmod{15625},\\
 (2195-\mathcal E)\cap(3296-\mathcal E)&=\{0\}
                         \pmod{15625}.
\end{align*}
The first supplies a clean clause band for every binary rotation; the
second supplies unique whole-cell anchor alignment after binary alignment.
The upper-dummy spill position $1064$ already belongs to $\mathcal E$.
Its unique possible odd spill cannot satisfy both distinct anchor tests.
These are the same synchronization arguments as the inherited exact-support
compiler. The support span $3296$ is less than half the circumference,
so no omitted circular case arises.

\subsection{Population padding can test guaranteed zero positions}

The two clause masks contribute $106$ set bits. Twenty overlap tests,
two anchor tests, and the low test contribute $23$ more. The remaining
$860$ bits needed to reach $989$ are supplied by $Z_{\rm pad}$.

For an aligned canonical witness, the largest radix exponent in any
$D_C$ convolution is $5495+3296=8791$; the corresponding bound for
$D_R$ is $42+3296=3338$, and temporal native support is at most $3296$.
All padding positions $9000,\ldots,9006$ are therefore identically zero.
Completeness loses no canonical witness. In soundness the added tests
may simply be discarded; adding zero tests cannot create a witness.

The exact population and low-bit invariants are
\begin{align*}
 \pc(M_F^{\rm native})&=106+23+860=989,\\
 \pc(M_C)&=d_{\rm base}-989,\\
 M_C&\equiv2\pmod4,&M_F^{\rm native}&\equiv4\pmod8,\\
 0<M_C,M_F^{\rm native}&<B-1.
\end{align*}
They preserve the inherited inverse-population theorem and its exact
two-bit surplus. The padding need not occur inside the clause band;
the proof uses its population and its vanishing on canonical outputs.

\subsection{Carries, margins, and five-adic tuning}

For the uniform inherited bound
\[
 \mathrm{mass}=(989+2)\left(2\sum_{e\in\mathcal E}c_e+6\right),
\]
$\mathrm{mass}$ has $118$ bits and $4\mathrm{mass}+8$ has $120$ bits.
The largest padding digit has $123$ bits, and
$2(2^{123}-1)+4<2^{125}$. Therefore $\rho=2^{125}$ exceeds every
required uniform carry threshold. Padding is part of the tested mask,
not an additional computation summand. All unshifted products and mask
positions remain inside the $15625$-cell circumference, and the inherited
strict-slack estimates remain available.

Because $125$ is a power of five, $\rho\equiv2\pmod5$.
The clause gap $1102\equiv2\pmod4$ makes paired reflected-clause
contributions cancel modulo five. The four bare shifts give
$D_C\equiv0\pmod5$, while $D_R\equiv4\pmod5$. Consequently
\[
 \boxed{2D_C-D_R\equiv1\pmod5.}
\]
The inherited dummy-control argument needs this quantity to be a unit;
its particular nonzero residue is immaterial. Independent power-of-five
width and height padding therefore still tunes the actual packed rank
to the needed temporal stride. This completes the coefficient-adapter
argument and the remaining part of Theorem~\ref{thm:main}.

\section{Materialized integers and operation-minimal alternatives}

The coefficient payloads are deterministic gzip files containing minimal
unsigned little-endian integers, with no header in the decompressed bytes.
The binding manifest records SHA-256 hashes of both forms. The sizes are:
\begin{center}\small
\begin{tabular}{@{}lrrr@{}}
\toprule
Port & Bits & Raw bytes & Gzip bytes\\
\midrule
$\beta$ & 1,953,125 & 244,141 & 273\\
$K_0$ & 1,958,376 & 244,797 & 10,170\\
$M_C$ & 1,953,125 & 244,141 & 373\\
$M_F$ & 1,953,126 & 244,141 & 371\\
\midrule
Total & & 977,220 & 11,187\\
\bottomrule
\end{tabular}
\end{center}
The old $943$-tile compact tuple occupies $5\,312\,269$ raw bytes,
and its largest coefficient has $13\,201\,251$ bits \cite{compact943}.
The new raw size is $81.604\%$ smaller, while the largest bit length
is $85.165\%$ smaller. These are actual coefficient-height improvements,
not merely better compression of unchanged integers.

\subsection{The unchanged 67-operation polynomial, with smaller coefficients}

For users prioritizing operation count, the archive includes the inherited
fifteen-positive-variable arithmetic DAG, rebound to~\eqref{eq:ports}.
Use the common interface with additional positive $f,\tau,y,\gamma$ and
without $r$. Set
\begin{align*}
 D_{\rm main}&=X+ac+\gamma H,&
 V&=c(Tf-1)-Rf^2,&K_{\rm aux}&=(i\Delta c^2)^2,\\
 N^{(0)}&=\tau^2-(XY^2)(XY^2+1)k^2,\\
 N^{(1)}&=D_{\rm main}^2-\Delta c^2,\\
 N^{(a)}&=K_{\rm aux}(V^2-y^2)+y^2,\\
 N^{(s)}&=\Delta f^2-K_{\rm aux}.
\end{align*}
Then
\[
 \mathcal P_{67}=N^{(0)}N^{(1)}N^{(a)}N_kN_tN^{(s)}-\Delta=0
\]
has the same positive existence predicate, degree $155$, and $67$
operations \cite{blank67}. Its arithmetic is not a new result here.
Its fixed coefficients are smaller by the new adapter.

The archive also includes package 22's unchanged seven-equation system
in sixteen positive unknowns, degree $24$ and $59$ operations, rebound
to the same tuple \cite{retained59}. These two baselines prevent a
twelve-variable improvement from being mistaken for an overall
operation-count record.

\section{Verification, reproducibility, and limits}
\label{sec:verification}

The Python programs use only the standard library. The delivered package
is designed to run after extraction, without relying on the source
repository layout, omitted historical files, or network access.
From the package root:
\begin{verbatim}
python3 reproduce.py
make pdf
\end{verbatim}
The first command regenerates the arithmetic circuits and coefficient
payloads, runs independent audits, and checks the fixed-instance bindings.
The second compiles this source using \texttt{pdflatex} twice.

\subsection{What the checks establish}

The arithmetic builder checks each of three variants against an independently
implemented eight-component conjugate-product evaluator on $24$ exact signed
assignments. It checks product and sum-of-squares finalizers, the parity
identity, dependency order, liveness, and operation counts. A separate
auditor imports no builder and checks the leading-form identities on $32$
exact assignments. The uniform proof in Section~\ref{sec:degree}, not the
finite sampling, establishes the degree for the actual large coefficients.

The coefficient verifier reads the actual bytes and independently
reconstructs the recipe. Its checks include $69\,230$ exact native-basis
identities, all $943$ valid tile clauses, rejection of $41\,492$ individual
copy/anchor mutations, rejection of all $44$ analogous empty-cell
mutations, all $944$ possible selector counts, support intersections,
population counts, low congruences, and padding support separation.
These finite basis identities establish the stated finite layout properties
exactly; the arbitrary-tableau transfer is the mathematical argument in
Section~\ref{sec:adapterproof}.

The release verifier additionally evaluates the supplied circuits modulo
three primes using residues of the actual four integers, and checks the
compact direct formula against the emitted circuit. These computations
audit binding and transcription. They are not a search for a full solution
or evidence concerning the truth of consistency.

\subsection{What is inherited and what remains open}

The new proof uses the same native compiler, Pell rank lemmas, half-binomial
criterion, and K215 calibration as the reviewed packages. Copies of the
critical inherited arguments and an immutable provenance ledger accompany
the archive. The new work has received a separate mathematical audit within
this assistant collaboration, not an external human endorsement.
No Lean or Rocq proof of the new norm or coefficient adapter is claimed.

The improvements do not settle any absolute minimization problem. The
nine-natural-variable form remains smaller in variable count; the
$67$-operation and $59$-operation baselines remain smaller in arithmetic
work. We have not constructed an eleven-positive-variable equation at
moderate degree or a new improvement to the signed $26$-integer family.

For contemporary context, Sun's September 2026 preprint states a constructive
ten-integer-variable representation for every recursively enumerable set
\cite{sun2026}. That general theorem gives a further route to a small
signed-variable count. This package does not instantiate its source-polynomial
compiler for K215 and does not present an unspecified polynomial as an
explicit ten-variable result. Likewise, the concrete Jones reduction used
by package 18 assumes a quartic input: feeding it a degree-$1081$ polynomial
without changing and proving the coding bounds would be invalid.

The most direct formalization project is now comparatively local: prove
Lemma~\ref{lem:gap}, the coherent-root recovery lemma, the scalar norm
identity, and the mixed-radix decoder, then compose them with an explicitly
formalized native compiler interface. The materialized tuple and full
circuits make this a fixed, reviewable target.

\appendix
\section{Package map and exact instance contract}

\begin{center}\small
\begin{tabular}{@{}p{.30\linewidth}p{.62\linewidth}@{}}
\toprule
Location & Contents\\
\midrule
\texttt{report.tex}, \texttt{report.pdf} & Complete article and compiled version.\\
\texttt{arithmetic/} & Self-contained emitter, nine circuit variants, algebra receipts.\\
\texttt{coefficients/} & Literal K215 table, ordered tiles, complete coefficient recipe,
actual compressed integers, independent byte verifier, binding manifest.\\
\texttt{baselines/} & The unchanged 67-operation polynomial and 59-operation system,
with complete new bindings.\\
\texttt{proofs/} & Expanded new proofs and independent mathematical reviews.\\
\texttt{upstream/} & Selected inherited proofs and endpoint source, preserved as evidence.\\
\texttt{verification/} & Independent arithmetic, fixed-instance, and release checks.\\
\texttt{provenance/} & Immutable source paths, revision, and byte hashes.\\
\texttt{licenses/} & Repository license and inherited machine notices.\\
\bottomrule
\end{tabular}
\end{center}

The canonical instance identifier is
\begin{center}\small\ttfamily
wang943-mixed125-external-padding-5pow9.
\end{center}
Circuit aliases $\mathtt{beta}$ and $\mathtt{Bm1}$ map to the payload
\texttt{beta}; $\mathtt{K0}$ and $\mathtt{Kconstant}$ map to
\texttt{K}; $\mathtt{MC}$ maps to \texttt{MC}; and $\mathtt{MF}$ maps
to \texttt{MF\_source}. All four ports must be replaced together.
The shifted mask convention is part of the theorem, not an optional
renaming. Historical source copies keep their historical bindings;
only the explicitly rebound delivered circuits select this instance.

\begin{thebibliography}{99}\small
\bibitem{catalogue} ProveIt contributors,
\emph{Con(ZFC) in Diophantine form: audited wave-3 comparison}.
\repo{SetTheory/Cardinals/docs/reports/hilbert-tenth-problem/con-zfc-diophantine-forms/README.md}{Pinned collection README}.

\bibitem{root20} ProveIt package 20,
\emph{Twelve positive unknowns by amalgamating the main Pell root}.
\repo{SetTheory/Cardinals/docs/reports/hilbert-tenth-problem/con-zfc-diophantine-forms/20-root12/amalgamation/ROOT12_PROOF.md}{Pinned proof}.

\bibitem{norm21} ProveIt package 21,
\emph{A twelve-positive-variable Diophantine equation for Con(ZFC)}.
\repo{SetTheory/Cardinals/docs/reports/hilbert-tenth-problem/con-zfc-diophantine-forms/21-norm12/proofs/NORM12_PROOF.md}{Pinned proof};
\repo{SetTheory/Cardinals/docs/reports/hilbert-tenth-problem/con-zfc-diophantine-forms/21-norm12/arithmetic/build_norm12.py}{scalar norm emitter}.

\bibitem{root09} ProveIt package 09, retained by package 20,
\emph{Root amalgamation and the signed-quotient interface}.
\repo{SetTheory/Cardinals/docs/reports/hilbert-tenth-problem/con-zfc-diophantine-forms/20-root12/amalgamation/dependencies/09_ROOT_AMALGAMATION_PROOF.md}{Pinned dependency}.

\bibitem{blank67} ProveIt, retained by package 21,
\emph{Independent proof and source audit: the 67-operation blank-input polynomial}.
\repo{SetTheory/Cardinals/docs/reports/hilbert-tenth-problem/con-zfc-diophantine-forms/21-norm12/upstream/FIXED_BLANK_67_PROOF.md}{Pinned proof}.

\bibitem{helical} ProveIt, retained by package 21,
\emph{Blank-input helical native compiler proof}.
\repo{SetTheory/Cardinals/docs/reports/hilbert-tenth-problem/con-zfc-diophantine-forms/21-norm12/upstream/BLANK_HELICAL_PROOF.md}{Pinned proof}.

\bibitem{wang943} ProveIt package 19,
\emph{Materialized four-coefficient Con(ZFC) instance: 943 signal tiles}.
\repo{SetTheory/Cardinals/docs/reports/hilbert-tenth-problem/con-zfc-diophantine-forms/19-wang943/compiler/WANG943_COEFFICIENTS.md}{Pinned language and compiler proof}.

\bibitem{compact943} ProveIt package 19,
\emph{Recommended compact instance: $B=2^{5^{10}}$}.
\repo{SetTheory/Cardinals/docs/reports/hilbert-tenth-problem/con-zfc-diophantine-forms/19-wang943/compiler/WANG943_COMPACT_LAYOUT.md}{Pinned exact-support layout}.

\bibitem{retained59} ProveIt package 22,
\emph{A 59-operation system in 16 positive unknowns}.
\repo{SetTheory/Cardinals/docs/reports/hilbert-tenth-problem/con-zfc-diophantine-forms/22-monic12/proofs/RETAINED59_PROOF.md}{Pinned proof}.

\bibitem{jones9} ProveIt package 18,
\emph{Three explicit Diophantine forms of ZFC consistency}.
\repo{SetTheory/Cardinals/docs/reports/hilbert-tenth-problem/con-zfc-diophantine-forms/18-jones9/report.tex}{Pinned report source};
\repo{SetTheory/Cardinals/docs/reports/hilbert-tenth-problem/con-zfc-diophantine-forms/18-jones9/nine/nine_section.tex}{concrete nine-variable section}.

\bibitem{kendpoint} ProveIt contributors,
\emph{KEndpoint215.blank\_halts\_iff\_zfc\_refutation}.
\repo{Computability/BusyBeaver/Lean/BusyBeaver/ZFC/KEndpoint215.lean}{Pinned Lean endpoint source}.

\bibitem{sun2026} Zhi-Wei Sun,
\emph{Ten unknowns for Hilbert's tenth problem over the integers},
arXiv:2609.01594v2, 2 September 2026.
\url{https://arxiv.org/html/2609.01594v2}.
\end{thebibliography}
\end{document}
